\documentclass[../DoAn.tex]{subfiles}
\begin{document}

\chapter{NGUỒN DỮ LIỆU}

\section*{A.1\quad Thống kê mẫu lành tính theo nguồn}
\addcontentsline{toc}{section}{A.1\quad Thống kê mẫu lành tính theo nguồn}

Bảng~\ref{tab:benign-sources} liệt kê 13~nguồn thu thập mẫu lành
tính, gồm bảy phiên bản Windows khác nhau từ ISO gốc của Microsoft
và sáu nguồn phần mềm hợp pháp bên ngoài.
\begin{table}[H]
\centering
\caption{Thống kê mẫu lành tính (benign) theo nguồn.}
\label{tab:benign-sources}
\begin{tabular}{|c|l|r|}
\hline
\multicolumn{1}{|c|}{\textbf{STT}} & \multicolumn{1}{c|}{\textbf{Nguồn}} & \multicolumn{1}{c|}{\textbf{Số mẫu}} \\
\hline
1  & Windows 10 LTSC 2021 (en-US x64)      & 18.009 \\
\hline
2  & Windows 11 Consumer 25H2 (en-US x64)  & 21.540 \\
\hline
3  & Windows Server 2022 (en-US x64)       & 15.801 \\
\hline
4  & Windows Server 2016 (en-US x64)       & 12.168 \\
\hline
5  & Windows Server 2019 (en-US x64)       & 12.875 \\
\hline
6  & Windows 8.1 Update 1 (en-US x64)      & 10.856 \\
\hline
7  & Windows 7 SP1 Ultimate (en-US x64)    &  8.821 \\
\hline
8  & Windows 10 22H2 (en-US x64)           &  5.275 \\
\hline
9  & Chocolatey packages (batch 1)         &  1.813 \\
\hline
10 & Chocolatey packages (batch 2)         &  1.091 \\
\hline
11 & Sysinternals Suite 2025-03            &    405 \\
\hline
12 & Nirsoft utilities collection          &    505 \\
\hline
13 & GitHub open-source releases           &    843 \\
\hline
   & \textbf{Tổng}                         & \textbf{110.002} \\
\hline
\end{tabular}
\end{table}
\section*{A.2\quad Thống kê mẫu mã độc theo nguồn}
\addcontentsline{toc}{section}{A.2\quad Thống kê mẫu mã độc theo nguồn}

Bảng~\ref{tab:malware-sources} liệt kê bốn tập VirusShare và
tập MalwareBazaar được thu thập làm mẫu mã độc.
\begin{table}[H]
\centering
\caption{Thống kê mẫu mã độc theo nguồn.}
\label{tab:malware-sources}
\begin{tabular}{|c|l|r|}
\hline
\multicolumn{1}{|c|}{\textbf{STT}} & \multicolumn{1}{c|}{\textbf{Nguồn}} & \multicolumn{1}{c|}{\textbf{Số mẫu}} \\
\hline
1 & VirusShare\_00499              & 31.042 \\
\hline
2 & VirusShare\_00495              & 32.891 \\
\hline
3 & VirusShare\_00483              & 38.536 \\
\hline
4 & VirusShare\_00480              & 28.603 \\
\hline
5 & MalwareBazaar (nhiều họ mã độc) &  1.404 \\
\hline
  & \textbf{Tổng}                  & \textbf{132.476} \\
\hline
\end{tabular}
\end{table}
\chapter{QUY TẮC ÁNH XẠ MITRE ATT\&CK VÀ THAM CHIẾU KỸ THUẬT}

\section*{B.1\quad Sysmon Event ID được theo dõi}
\addcontentsline{toc}{section}{B.1\quad Sysmon Event ID được theo dõi}

Bảng~\ref{tab:sysmon-eids} liệt kê toàn bộ 23~Event ID Sysmon
được ghi nhận trong phiên phân tích cùng nhóm hành vi tương ứng.
\begin{table}[H]
\centering
\caption{Các Sysmon Event ID được theo dõi và nhóm hành vi tương ứng.}
\label{tab:sysmon-eids}
\begin{tabular}{|l|l|l|}
\hline
\multicolumn{1}{|c|}{\textbf{Event ID}} & \multicolumn{1}{c|}{\textbf{Nhóm hành vi}} & \multicolumn{1}{c|}{\textbf{Mô tả}} \\
\hline
1          & Tiến trình  & Tạo tiến trình \\
\hline
2          & Tệp         & Thay đổi thời gian tạo tệp \\
\hline
3          & Mạng        & Kết nối mạng \\
\hline
5          & Tiến trình  & Kết thúc tiến trình \\
\hline
6          & Khác        & Nạp driver \\
\hline
7          & DLL         & Nạp thư viện động \\
\hline
8          & Tiêm mã     & Tạo luồng từ xa (remote thread) \\
\hline
9          & Khác        & Đọc thô đĩa cứng \\
\hline
10         & Tiêm mã     & Truy cập tiến trình khác \\
\hline
11         & Tệp         & Tạo tệp \\
\hline
12, 13, 14 & Registry    & Tạo/xóa khóa, ghi giá trị, đổi tên khóa registry \\
\hline
15         & Tệp         & Tạo luồng dữ liệu thay thế (ADS) \\
\hline
17, 18     & Named pipe  & Tạo / kết nối named pipe \\
\hline
19, 20, 21 & WMI         & Đăng ký WMI filter, consumer, binding \\
\hline
22         & DNS         & Truy vấn DNS \\
\hline
23, 26     & Tệp         & Xóa tệp \\
\hline
25         & Tiêm mã     & Giả mạo tiến trình \\
\hline
\end{tabular}
\end{table}
\section*{B.2\quad Các luật ánh xạ MITRE ATT\&CK đã triển khai}
\addcontentsline{toc}{section}{B.2\quad Các luật ánh xạ MITRE ATT\&CK đã triển khai}

Bảng~\ref{tab:mitre-rules} liệt kê toàn bộ 23~kỹ thuật ánh xạ MITRE
ATT\&CK được triển khai trong bộ phân tích hành vi. Bốn kỹ thuật
T1071.001, T1566.001, T1070.007 và T1568.002 được đánh mức tin cậy
trung bình do tín hiệu kích hoạt có thể xuất hiện trong hoạt động
bình thường; các kỹ thuật còn lại được đánh mức cao.
\begin{table}[H]
\centering
\caption{Các luật ánh xạ MITRE ATT\&CK được triển khai trong bộ
  phân tích hành vi.}
\label{tab:mitre-rules}
\begin{tabular}{|l|l|>{\raggedright\arraybackslash}p{5.8cm}|}
\hline
\multicolumn{1}{|c|}{\textbf{Kỹ thuật}} & \multicolumn{1}{c|}{\textbf{Tên}} & \multicolumn{1}{c|}{\textbf{Tín hiệu kích hoạt}} \\
\hline
T1003.001   & LSASS Memory                & Process access nhắm vào \texttt{lsass.exe} \\
\hline
T1053.005   & Scheduled Task              & Ghi vào khóa \texttt{ScheduledTask} \\
\hline
T1055       & Process Injection           & Remote thread (EID~8) hoặc process access (EID~10) \\
\hline
T1055.012   & Process Hollowing           & Giả mạo tiến trình (EID~25) \\
\hline
T1070       & Indicator Removal           & Xóa tệp log, shadow, backup hoặc .evt \\
\hline
T1070.007   & Clear Persistence           & Xóa khóa registry \\
\hline
T1071.001   & App.\ Layer Protocol        & Kết nối ra IP công khai ngoài mạng nội bộ \\
\hline
T1071.003   & Named Pipe C2               & Named pipe khớp từ khóa C2 đã biết \\
\hline
T1090.003   & Multi-hop Proxy             & Truy vấn DNS tên miền \texttt{.onion} \\
\hline
T1218       & Signed Binary Proxy Exec.   & Tiến trình là công cụ LOLBin \\
\hline
T1490       & Inhibit System Recovery     & Thao tác tệp shadow, VSS hoặc backup \\
\hline
T1543.003   & Windows Service             & Ghi vào khóa \texttt{Services} \\
\hline
T1546.003   & WMI Event Subscription      & Đăng ký WMI filter, consumer, binding \\
\hline
T1546.010   & AppInit DLLs                & Ghi vào khóa \texttt{AppInit\_DLLs} \\
\hline
T1546.012   & IFEO Injection              & Ghi vào \texttt{Image File Execution} \\
\hline
T1547.001   & Registry Run Keys           & Khóa \texttt{CurrentVersion\textbackslash{}Run} hoặc Startup \\
\hline
T1547.004   & Winlogon Helper DLL         & Ghi vào khóa \texttt{Winlogon} \\
\hline
T1547.005   & Security Support Provider   & Ghi vào khóa \texttt{SecurityProviders} \\
\hline
T1547.008   & LSASS Driver                & Ghi vào \texttt{Lsa\textbackslash{}Notification} \\
\hline
T1566.001   & Spearphishing Attachment    & Tiến trình cha là ứng dụng Office hoặc Acrobat \\
\hline
T1568.002   & Domain Generation Alg.      & Entropy Shannon nhãn phụ $>3{,}5$~bit và độ dài $\geq 12$~ký tự \\
\hline
T1571       & Non-Standard Port           & Kết nối đến cổng C2 (4444, 1337, 8888, 9999, 31337, 8080, 8443) \\
\hline
T1620       & Reflective Code Loading     & Assembly .NET nạp qua \texttt{System.Reflection} \\
\hline
\end{tabular}
\end{table}
\section*{B.3\quad Endpoint và sự kiện SSE của server web demo}
\addcontentsline{toc}{section}{B.3\quad Endpoint và sự kiện SSE của server web demo}

Bảng~\ref{tab:demo-endpoints} liệt kê tám endpoint REST của server;
Bảng~\ref{tab:sse-events} mô tả chín loại sự kiện trong luồng SSE.
\begin{table}[H]
\centering
\caption{Toàn bộ endpoint của server web demo.}
\label{tab:demo-endpoints}
\begin{tabular}{|l|l|>{\raggedright\arraybackslash}p{6.5cm}|}
\hline
\multicolumn{1}{|c|}{\textbf{Endpoint}} & \multicolumn{1}{c|}{\textbf{Phương thức}} & \multicolumn{1}{c|}{\textbf{Chức năng}} \\
\hline
\texttt{/}                          & GET  & Phục vụ trang giao diện (\texttt{index.html}) \\
\hline
\texttt{/api/analyze}               & POST & Nhận tệp mẫu và tham số; tạo \texttt{job\_id}, lưu vào \texttt{uploads/}, khởi chạy pipeline trong luồng nền \\
\hline
\texttt{/api/jobs/<id>/stream}      & GET  & Luồng SSE; phát sự kiện pipeline theo thời gian thực \\
\hline
\texttt{/api/jobs/<id>/stop}        & POST & Đặt cờ dừng, đẩy \texttt{stopped}+\texttt{eof} vào hàng đợi, tắt VM qua \texttt{VMManager} \\
\hline
\texttt{/api/jobs/<id>}             & GET  & Trả về trạng thái và kết quả cuối của một job \\
\hline
\texttt{/api/reports}               & GET  & Liệt kê tệp JSON trong \texttt{results/pipeline\_runs/}, sắp xếp theo thời gian giảm dần \\
\hline
\texttt{/api/report/<filename>}     & GET    & Đọc và trả về nội dung một báo cáo JSON \\
\hline
\texttt{/api/report/<filename>}     & DELETE & Xóa tệp báo cáo JSON khỏi thư mục \texttt{results/pipeline\_runs/} \\
\hline
\end{tabular}
\end{table}
\begin{table}[H]
\centering
\caption{Chín loại sự kiện SSE trong luồng \texttt{/api/jobs/<id>/stream}.}
\label{tab:sse-events}
\begin{tabular}{|l|>{\raggedright\arraybackslash}p{9.5cm}|}
\hline
\multicolumn{1}{|c|}{\textbf{Loại sự kiện}} & \multicolumn{1}{c|}{\textbf{Nội dung}} \\
\hline
\texttt{ping}         & Giữ kết nối; client bỏ qua \\
\hline
\texttt{started}      & Pipeline bắt đầu; kèm tên tệp mẫu \\
\hline
\texttt{stage\_start} & Giai đoạn bắt đầu; kèm số thứ tự \texttt{n} (1–12) và tên \\
\hline
\texttt{stage\_done}  & Giai đoạn hoàn thành; kèm số thứ tự \texttt{n} \\
\hline
\texttt{log}          & Dòng nhật ký; kèm cấp độ (INFO/WARNING/ERROR) và nội dung \\
\hline
\texttt{done}         & Pipeline hoàn tất; kèm toàn bộ báo cáo JSON \\
\hline
\texttt{stopped}      & Pipeline bị người dùng dừng \\
\hline
\texttt{error}        & Pipeline gặp lỗi; kèm thông báo lỗi \\
\hline
\texttt{eof}          & Kết thúc luồng; client đóng \texttt{EventSource} \\
\hline
\end{tabular}
\end{table}
\section*{B.4\quad Danh sách lọc nhiễu trong giai đoạn động}
\addcontentsline{toc}{section}{B.4\quad Danh sách lọc nhiễu trong giai đoạn động}

Mục này tổng hợp các danh sách lọc nhiễu dùng ở giai đoạn động.
Bảng~\ref{tab:silketw-exclude} liệt kê tám tiến trình thuộc
\texttt{\_DOTNET\_NOISE\_PROCS} bị loại khỏi đầu ra SilkETW;
Bảng~\ref{tab:silketw-keywords} liệt kê bảy từ khóa namespace
thuộc \texttt{\_SUSPICIOUS\_NS} dùng để đánh dấu nghi ngờ;
Bảng~\ref{tab:tree-noise} liệt kê 27~tiến trình hệ thống thường
trực bị lọc khỏi cây tiến trình (\texttt{\_TREE\_NOISE}).
\begin{table}[H]
\centering
\caption{Tám tiến trình trong \texttt{\_DOTNET\_NOISE\_PROCS} bị loại
  khỏi đầu ra bộ chuẩn hóa SilkETW.}
\label{tab:silketw-exclude}
\begin{tabular}{|>{\raggedright\arraybackslash}p{4.2cm}|>{\raggedright\arraybackslash}p{8cm}|}
\hline
\multicolumn{1}{|c|}{\textbf{Tiến trình}} & \multicolumn{1}{c|}{\textbf{Lý do loại trừ}} \\
\hline
\texttt{msbuild}, \texttt{devenv}, \texttt{dotnet},
\texttt{vstest}, \texttt{testhost} & Công cụ phát triển .NET, không phải đối tượng phân tích \\
\hline
\texttt{silketw} & Tránh tự tham chiếu của công cụ thu thập \\
\hline
\texttt{powershell}, \texttt{pwsh} & Sysmon EID~1 đã ghi đầy đủ khi mẫu gọi shell \\
\hline
\end{tabular}
\end{table}
\begin{table}[H]
\centering
\caption{Bảy từ khóa trong \texttt{\_SUSPICIOUS\_NS} dùng đánh dấu
  sự kiện SilkETW nghi ngờ.}
\label{tab:silketw-keywords}
\begin{tabular}{|l|>{\raggedright\arraybackslash}p{8.5cm}|}
\hline
\multicolumn{1}{|c|}{\textbf{Từ khóa namespace}} & \multicolumn{1}{c|}{\textbf{Hành vi phản ánh}} \\
\hline
\texttt{dynamicClass}      & Tạo kiểu .NET động tại thời điểm chạy \\
\hline
\texttt{Invoke}            & Gọi phương thức qua Reflection \\
\hline
\texttt{Emit}              & Tự sinh mã IL lúc chạy (IL Emit) \\
\hline
\texttt{Reflection}        & Truy cập metadata và bộ nhớ qua introspection \\
\hline
\texttt{RunspaceFactory}   & Tạo môi trường PowerShell trong bộ nhớ \\
\hline
\texttt{AmsiUtils}         & Vô hiệu hóa Antimalware Scan Interface \\
\hline
\texttt{System.Net}        & Kết nối mạng từ bên trong .NET Runtime \\
\hline
\end{tabular}
\end{table}

27~tiến trình hệ thống thường trực sau đây bị lọc khỏi cây tiến
trình (hằng \texttt{\_TREE\_NOISE}), chia theo bốn nhóm; tiến trình
chỉ bị ẩn khi toàn bộ nhánh con cháu của nó cũng là nhiễu. Tên đã
bỏ hậu tố \texttt{.exe}.
\begin{table}[H]
\centering
\caption{27 tiến trình hệ thống bị lọc khỏi cây tiến trình
  (hằng \texttt{\_TREE\_NOISE}).}
\label{tab:tree-noise}
\begin{tabular}{|l|>{\raggedright\arraybackslash}p{8.5cm}|}
\hline
\multicolumn{1}{|c|}{\textbf{Nhóm}} & \multicolumn{1}{c|}{\textbf{Tiến trình}} \\
\hline
Lõi và phiên Windows &
  \texttt{csrss}, \texttt{smss}, \texttt{wininit}, \texttt{winlogon},
  \texttt{services}, \texttt{lsm}, \texttt{lsass}, \texttt{svchost},
  \texttt{dllhost}, \texttt{conhost}, \texttt{taskhostw},
  \texttt{runtimebroker}, \texttt{sihost}, \texttt{ctfmon},
  \texttt{fontdrvhost}, \texttt{dwm} \\
\hline
Dịch vụ hệ thống và WMI &
  \texttt{spoolsv}, \texttt{msdtc}, \texttt{audiodg},
  \texttt{wmiprvse}, \texttt{wsmprovhost} \\
\hline
Tìm kiếm và bảo mật Windows &
  \texttt{searchindexer}, \texttt{searchprotocolhost},
  \texttt{securityhealthservice}, \texttt{securityhealthsystray} \\
\hline
Tiến trình VirtualBox &
  \texttt{vboxservice}, \texttt{vboxtray} \\
\hline
\end{tabular}
\end{table}

\section*{B.5\quad Endpoint của Agent HTTP trong máy ảo}
\addcontentsline{toc}{section}{B.5\quad Endpoint của Agent HTTP trong máy ảo}

Bảng~\ref{tab:agent-endpoints} liệt kê toàn bộ 11~endpoint của
Agent HTTP chạy trong máy ảo Windows.
\begin{table}[H]
\centering
\caption{Toàn bộ endpoint của Agent HTTP trong máy ảo.}
\label{tab:agent-endpoints}
\begin{tabular}{|l|l|>{\raggedright\arraybackslash}p{6.5cm}|}
\hline
\multicolumn{1}{|c|}{\textbf{Endpoint}} & \multicolumn{1}{c|}{\textbf{Phương thức}} & \multicolumn{1}{c|}{\textbf{Chức năng}} \\
\hline
\texttt{/health}                 & GET  & Kiểm tra Agent đã khởi động; trả về trạng thái và thời gian hoạt động \\
\hline
\texttt{/execute}                & POST & Khởi chạy tệp mẫu dưới dạng tiến trình mới; trả về PID \\
\hline
\texttt{/shell}                  & POST & Thực thi lệnh shell tùy ý, đồng bộ hoặc nền (\texttt{async=true}) \\
\hline
\texttt{/upload}                 & POST & Ghi nội dung văn bản vào đường dẫn tùy chỉnh trên guest \\
\hline
\texttt{/read}                   & GET  & Đọc nội dung tệp từ hệ thống tệp guest; Bridge dùng để poll log Sysmon \\
\hline
\texttt{/process\_tree/\{pid\}}  & GET  & Liệt kê toàn bộ PID con và cháu của tiến trình gốc \\
\hline
\texttt{/process/\{pid\}}        & GET  & Kiểm tra trạng thái và thời gian chạy của một PID cụ thể \\
\hline
\texttt{/status}                 & GET  & Trạng thái tổng thể tiến trình hiện tại; tự dừng nếu vượt timeout \\
\hline
\texttt{/stop}                   & POST & Dừng tiến trình đang chạy theo PID \\
\hline
\texttt{/logs}                   & GET  & Trả về nhật ký nội bộ của Agent (sự kiện execute, stop, \ldots) \\
\hline
\texttt{/screenshot}             & GET  & Chụp màn hình guest, trả về ảnh dạng Base64 PNG \\
\hline
\end{tabular}
\end{table}
\section*{B.6\quad Ánh xạ sự kiện chuẩn hóa vào nhóm hành vi}
\addcontentsline{toc}{section}{B.6\quad Ánh xạ sự kiện chuẩn hóa vào nhóm hành vi}

Bảng~\ref{tab:event-groups} liệt kê 11~nhóm hành vi và các loại
sự kiện chuẩn hóa tương ứng dùng trong mô-đun xây dựng chuỗi hành vi.
\begin{table}[H]
\centering
\caption{Ánh xạ 24~loại sự kiện chuẩn hóa vào 11~nhóm hành vi.}
\label{tab:event-groups}
\begin{tabular}{|l|>{\raggedright\arraybackslash}p{9cm}|}
\hline
\multicolumn{1}{|c|}{\textbf{Nhóm hành vi}} & \multicolumn{1}{c|}{\textbf{Loại sự kiện}} \\
\hline
Tiến trình         & Tạo tiến trình, kết thúc tiến trình \\
\hline
Tiêm mã            & Truy cập tiến trình khác, tạo luồng từ xa, giả mạo tiến trình \\
\hline
Nạp thư viện động & Nạp DLL \\
\hline
Tệp               & Tạo tệp, xóa tệp, xóa tệp có ghi nhận, thay đổi thời gian tệp, tạo ADS \\
\hline
Registry           & Tạo khóa, ghi khóa, xóa khóa \\
\hline
Mạng               & Kết nối mạng \\
\hline
DNS                & Truy vấn DNS \\
\hline
Named pipe         & Tạo pipe, kết nối pipe \\
\hline
WMI                & Đăng ký WMI filter, consumer, binding \\
\hline
.NET Runtime       & Nạp assembly .NET \\
\hline
Khác               & Nạp driver, đọc đĩa thô \\
\hline
\end{tabular}
\end{table}
\section*{B.7\quad Cấu trúc báo cáo JSON đầu ra của pipeline}
\addcontentsline{toc}{section}{B.7\quad Cấu trúc báo cáo JSON đầu ra của pipeline}

Bảng~\ref{tab:report-impl} liệt kê tên trường thực tế của từng
nhóm trong báo cáo JSON theo thiết kế tại Mục~\ref{section:3.5.2}.
\begin{table}[H]
\centering
\caption{Cấu trúc báo cáo JSON đầu ra của pipeline.}
\label{tab:report-impl}
\begin{tabular}{|l|>{\raggedright\arraybackslash}p{9cm}|}
\hline
\multicolumn{1}{|c|}{\textbf{Nhóm}} & \multicolumn{1}{c|}{\textbf{Trường JSON}} \\
\hline
Kết quả phân loại   &
  \texttt{verdict}, \texttt{confidence},
  \texttt{ml\_label}, \texttt{ml\_prob},
  \texttt{static\_result} (xác suất ML thô và thông tin PE cơ bản) \\
\hline
Thông tin hành vi   &
  \texttt{behavioral\_chain}, \texttt{process\_tree},
  \texttt{narrative}, \texttt{mitre\_techniques},
  \texttt{mitre\_details} (mỗi kỹ thuật kèm tên, nguồn và mức tin cậy) \\
\hline
Chỉ số xâm phạm    &
  \texttt{ioc} (10 danh mục: ips, domains, urls,
  files\_created, files\_deleted, registry\_keys,
  commands, named\_pipes, injected\_pids, hashes) \\
\hline
Thông tin phiên     &
  \texttt{metadata}: sha256, file\_size, started\_at,
  ended\_at, duration\_s, stages, event\_count,
  ttp\_count, analyzed\_at \\
\hline
Cảnh báo và lỗi    &
  \texttt{warnings}, \texttt{errors} \\
\hline
\end{tabular}
\end{table}
\end{document}
